\BeleskeID{09_2-s032}
% element: 09_2-s032:e04
Iz uslova se vidi da postoji oblast za pretpostavljeni rad tranzistora. Za P koristi se struja zasićenja, a za N omski izraz:
\[\begin{aligned}&\frac{k_p}2\frac1{1+\frac{V_{GS,P}-V_{Tp}}{L_pE_{Cp}}}(V_{GS,P}-V_{Tp})^2\\&\quad=\frac{k_n}2\frac1{1+\frac{V_{DS,N}}{L_nE_{Cn}}}\left(2V_{DS,N}(V_{GS,N}-V_{Tn})-V_{DS,N}^2\right)\end{aligned}\]
Zamenjuju se $V_{GS,P}=V_I-V_{DD}$, $V_{GS,N}=V_I$ i $V_{DS,N}=V_O$. Jednačina ravnoteže struja zatim povezuje ulazni i izlazni napon u četvrtoj oblasti karakteristike.
